%
%   ABSTRACTS
%

\pdfbookmark[0]{Abstracts}{Abstracts}



%%%%%%%%%% Castellano %%%%%%%%%%%%
                                                                   
\pdfbookmark[1]{Resumen (Spanish abstract)}{Resumo (Spanish abstract)}
\chapter*{Resumen}
\markboth{\color[cmyk]{.0, 1.00, .37, .19} Resumen (thesis abstract in Spanish)}{}
\markright{\color[cmyk]{.0, 1.00, .37, .19} Resumen (thesis abstract in Spanish)}{}
\label{sec:abstract_spanish}

El estudio del microbioma humano ha avanzado de forma acelerada gracias a la caída del coste de la secuenciación y a la disponibilidad de grandes repositorios públicos. Este progreso ha confirmado asociaciones entre perfiles microbianos y múltiples patologías, pero también ha puesto de relieve retos persistentes: la naturaleza composicional de los datos procedentes del microbioma, el ruido técnico y la variabilidad interindividual. En este contexto, los enfoques de \textit{Machine Learning} brindan una ventaja sustancial para detectar señales complejas con posible significado biológico y estimar con rigor el valor predictivo de los perfiles microbianos en entornos clínicos. 

Esta Tesis Doctoral responde a dichos desafíos adoptando un enfoque computacional estandarizado y reproducible, desde el procesado de lecturas 16S rRNA hasta la aplicación de modelos de \textit{Machine Learning} para el análisis de estos datos. Bajo este enfoque computacional, se identifican patrones microbianos que distinguen de forma consistente distintos estados clínicos en enfermedades metabólicas (síndrome metabólico y diabetes tipo 2); se obtienen modelos predictivos para diabetes tipo 1 sustentados en un conjunto acotado de taxones con alta relevancia biológica; y, en vaginosis bacteriana, se propone un subtipado alternativo del microbioma vaginal que añade información complementaria no recogida por otros métodos de estratificación, y puede apoyar la predicción de la respuesta terapéutica. Estos resultados se presentan en un compendio de tres publicaciones científicas. 

En conjunto, esta Tesis Doctoral ofrece evidencia de que los perfiles microbianos permiten la caracterización de estados clínicos y ofrece un marco metodológico reproducible que facilita su evaluación comparativa entre cohortes. 


%%%%%%%%%% Galego %%%%%%%%%%%%


\pdfbookmark[1]{Resumo (Galician abstract)}{Resumo (Galician abstract)}
\chapter*{Resumo}
\markboth{\color[cmyk]{.0, 1.00, .37, .19} Resumo (thesis abstract in Galician)}{}
\markright{\color[cmyk]{.0, 1.00, .37, .19} Resumo (thesis abstract in Galician)}{}
\label{sec:abstract_galician}

O estudo do microbioma humano avanzou de forma acelerada grazas á caída do custo da secuenciación e á dispoñibilidade de grandes repositorios públicos. Este progreso confirmou asociacións entre perfís microbianos e múltiples patoloxías, pero tamén puxo de relevo retos persistentes: a natureza composicional dos datos procedentes do microbioma, o ruído técnico e a variabilidade interindividual. Neste contexto, os enfoques de \textit{Machine Learning} brindan unha vantaxe substancial para detectar sinais complexos con posible significado biolóxico e estimar con rigor o valor preditivo dos perfís microbianos en escenarios clínicos.

Esta Tese Doutoral responde aos devanditos desafíos adoptando un enfoque computacional estandarizado e reproducible, desde o procesado de lecturas 16S rRNA ata a aplicación de modelos de \textit{Machine Learning} para a análise destes datos. Baixo este enfoque computacional, identifícanse patróns microbianos que distinguen de forma consistente distintos estados clínicos en enfermidades metabólicas (síndrome metabólico e diabetes tipo 2); obtéñense modelos preditivos para diabetes tipo 1 sustentados nun conxunto acoutado de taxóns con alta relevancia biolóxica; e, en vaxinose bacteriana, proponse un subtipado alternativo do microbioma vaxinal que engade información complementaria non recollida por outros métodos de estratificación, e pode apoiar a predición da resposta terapéutica. Estes resultados preséntanse nun compendio de tres publicacións científicas.

En conxunto, esta Tese Doutoral ofrece evidencia de que os perfís microbianos permiten a caracterización de estados clínicos e ofrece un marco metodolóxico reproducible que facilita a súa avaliación comparativa entre cohortes.